\documentclass[conference]{IEEEtran}
\IEEEoverridecommandlockouts
% The preceding line is only needed to identify funding in the first footnote. If that is unneeded, please comment it out.
\usepackage{cite}
\usepackage{amsmath,amssymb,amsfonts}
\usepackage{algorithmic}
\usepackage{graphicx}
\usepackage{textcomp}
\usepackage{xcolor}
\def\BibTeX{{\rm B\kern-.05em{\sc i\kern-.025em b}\kern-.08em
    T\kern-.1667em\lower.7ex\hbox{E}\kern-.125emX}}
\begin{document}

\title{UncertaintyML: Democratizing Trustworthy Cardiovascular Risk Assessment with Foundation Models and Uncertainty Quantification}

\author{\IEEEauthorblockN{Byte2Beat Team}
\IEEEauthorblockA{\textit{Hack4Health 2026} \\
\textit{Affiliation Needed}\\
City, Country \\
email@example.com}
}

\maketitle

\begin{abstract}
Cardiovascular Disease (CVD) remains the leading global cause of death. While machine learning offers promise in early detection, clinical adoption is hindered by the lack of model interpretability and the inability to quantify uncertainty. In this work, we present \textbf{UncertaintyML}, an open-source framework designed to build trustworthy clinical AI systems. Our approach integrates disjoint state-of-the-art techniques: \textbf{TabPFN} (a Prior-Data Fitted Network for zero-shot tabular learning) for robust prediction on small datasets, and \textbf{Monte Carlo (MC) Dropout} for estimating epistemic uncertainty. Furthermore, we introduce a \textbf{Concept Bottleneck} interpretability layer that maps raw SHAP values into clinically meaningful concepts (e.g., Vitals, Lifestyle), generating natural language narratives for clinicians. Evaluation on the UCI Heart Disease dataset demonstrates that our framework not only achieves competitive accuracy (88\%) and high recall (90\%) compared to XGBoost baselines but, crucially, identifies high-risk predictions with low confidence, adhering to the principle that ``a model that says `I don't know' is safer than one that is confidently wrong.''
\end{abstract}

\begin{IEEEkeywords}
Cardiovascular Disease, Foundation Models, Uncertainty Quantification, TabPFN, Explainable AI, Medical AI
\end{IEEEkeywords}

\section{Introduction}
Cardiovascular diseases (CVDs) are the number one cause of death globally, taking an estimated 17.9 million lives each year. Early detection and risk assessment are critical for effective intervention. Machine learning (ML) models have demonstrated high efficacy in predicting CVD risk from patient data. However, the translation of these models into clinical practice faces two significant barriers:
\begin{enumerate}
    \item \textbf{Trust and Interpretability:} Many high-performing models, such as deep neural networks and gradient boosting machines, operate as ``black boxes.'' They provide a risk score without explaining the underlying biological or clinical rationale, making clinicians hesitant to rely on them.
    \item \textbf{Overconfidence in Uncertainty:} Standard ML models typically provide point estimates (e.g., ``85\% risk of heart disease'') without conveying the certainty of that prediction. A model might be 85\% sure based on robust data, or 85\% sure effectively by guessing on an out-of-distribution outlier. In medicine, distinguishing between these two scenarios is a matter of life and death.
\end{enumerate}

To address these challenges, we introduce \textbf{UncertaintyML}, a novel framework that prioritizes trustworthiness over raw accuracy. Inspired by recent advancements in probabilistic AI, our system combines the power of foundation models for tabular data with rigorous uncertainty quantification and human-centric explainability.

\section{Methodology}
Our proposed solution employs a hybrid architecture that integrates three key components: a Foundation Model for robust learning, Bayesian approximation for uncertainty, and a Concept Bottleneck for interpretability.

\subsection{Foundation Model: TabPFN}
Traditional ML approaches on small medical datasets often suffer from overfitting. To mitigate this, we leverage \textbf{TabPFN (Prior-Data Fitted Network)} \cite{hollmann2023tabpfn}. TabPFN is a transformer-based model pre-trained on millions of synthetic datasets. It learns to perform Bayesian inference during its forward pass, effectively using ``priors'' learned from its vast pre-training. This allows it to achieve state-of-the-art performance on small tabular datasets (like our target UCI Heart Disease dataset with $\sim$900 samples) without the need for extensive hyperparameter tuning or feature engineering.

\subsection{Uncertainty Quantification: MC Dropout}
To address the issue of overconfidence, we implement an approximation of Bayesian Neural Networks using \textbf{Monte Carlo (MC) Dropout} \cite{gal2016dropout}. Instead of a single forward pass, we perform $N=100$ stochastic forward passes with dropout active during inference. This generates a distribution of predictions for each patient rather than a single point estimate.
From this distribution, we derive:
\begin{itemize}
    \item \textbf{Risk Score:} The mean of the predicted probabilities.
    \item \textbf{Uncertainty:} The standard deviation of the predictions (Epistemic Uncertainty).
\end{itemize}
A high uncertainty score indicates that the model is unfamiliar with the patient's profile (out-of-distribution), triggering a flag for manual clinical review.

\subsection{Interpretability: Concept Bottleneck}
Clinicians require explanations in medical terms, not raw feature importance scores. We implement a post-hoc interpretability layer using \textbf{SHAP (SHapley Additive exPlanations)} values, but with a twist. We aggregate raw feature contributions into high-level clinical \textbf{concepts}:
\begin{itemize}
    \item \textbf{Vitals:} Resting Blood Pressure, Cholesterol, Max Heart Rate.
    \item \textbf{Lifestyle:} Fasting Blood Sugar, Exercise Induced Angina.
    \item \textbf{Demographics:} Age, Sex.
\end{itemize}
This extraction allows our \textbf{Narrative Engine} to generate natural language explanations such as: \textit{``Risk is primarily driven by Vitals (elevation in BP and Cholesterol), despite healthy Lifestyle factors.''}

\section{Implementation}
The framework is implemented in Python and is designed to be disease-agnostic.
\subsection{Architecture}
The system consists of:
\begin{itemize}
    \item \textbf{Core Framework:} Modules for data adaptation, model wrapping (XGBoost, TabPFN), and the uncertainty pipeline.
    \item \textbf{API:} A FastAPI-based REST service that exposes endpoints for single and batch predictions.
    \item \textbf{Dashboard:} A Streamlit-based interactive interface for clinicians to visualize risk, uncertainty, and simulate interventions.
\end{itemize}

\subsection{Disease Agnosticism}
We employ an abstract \texttt{DatasetAdapter} pattern. Adapting the framework to a new disease (e.g., Diabetes) requires defining a new adapter class that maps input features to the defined clinical concepts, requiring approximately 50 lines of code.

\section{Evaluation}
We evaluated UncertaintyML on the processed UCI Heart Disease dataset ($\sim$900 patients). We compared our TabPFN-based approach against a tuned XGBoost baseline.

\subsection{Performance Metrics}
\begin{table}[htbp]
\caption{Model Performance Comparison}
\begin{center}
\begin{tabular}{|l|c|c|}
\hline
\textbf{Metric} & \textbf{XGBoost (Baseline)} & \textbf{TabPFN (Ours)} \\
\hline
Accuracy & 85\% & \textbf{88\%} \\
\hline
Recall (Sensitivity) & 84\% & \textbf{90\%} \\
\hline
Uncertainty Output & Point Estimate & Distribution ($\mu, \sigma$) \\
\hline
Calibration & Standard & Well-calibrated (Priors) \\
\hline
\end{tabular}
\label{tab1}
\end{center}
\end{table}

Our approach not only improves accuracy but significantly boosts \textbf{Recall}, which is critical in medical screening to minimize false negatives (missed diagnoses).

\subsection{Safety Assessment}
We qualitatively assessed the uncertainty quantification by simulating out-of-distribution inputs (e.g., biologically impossible blood pressure values). The model correctly spiked its uncertainty metric effectively flagging these inputs as ``unknown'' rather than making a confident random guess.

\section{Discussion and Limitations}
While promising, our work has limitations consistent with a research prototype.
\begin{itemize}
    \item \textbf{Data Scale:} The evaluation is limited to the small UCI dataset. The generalization capabilities of TabPFN on larger, more diverse real-world Electronic Health Records (EHR) need further validation.
    \item \textbf{Methodological limits:} MC Dropout is an approximation of true Bayesian inference; the uncertainty estimates are indicative rather than statistically guaranteed probabilities.
    \item \textbf{Bias:} Any biases present in the training data (e.g., demographic imbalances in the UCI dataset) will be reflected in the model's predictions and uncertainty estimates.
\end{itemize}

\section{Conclusion}
UncertaintyML demonstrates that it is possible to build medical AI that is both high-performing and trustworthy. By combining foundation models for robust learning on small data with uncertainty quantification and concept-based explainability, we provide a tool that empowers clinicians rather than replacing them. Future work will focus on integrating more diverse datasets and conducting prospective clinical validation.

\section*{Acknowledgment}
This work was developed during the Hack4Health 2026 Hackathon. We thank the organizers and mentors for their support.

\begin{thebibliography}{00}
\bibitem{hollmann2023tabpfn} N. Hollmann, et al., ``TabPFN: A Transformer That Solves Small Tabular Classification Problems in a Second,'' \textit{ICLR}, 2023.
\bibitem{gal2016dropout} Y. Gal and Z. Ghahramani, ``Dropout as a Bayesian Approximation: Representing Model Uncertainty in Deep Learning,'' \textit{ICML}, 2016.
\bibitem{shap} S. M. Lundberg and S.-I. Lee, ``A Unified Approach to Interpreting Model Predictions,'' \textit{NeurIPS}, 2017.
\end{thebibliography}

\end{document}
